% =============================================================================
%  ONE-FILE LaTeX CV TEMPLATE (beginner friendly)
% =============================================================================
%
%  WHAT IS THIS?
%  LaTeX is a typesetting language: you write plain text plus a few commands,
%  and a compiler turns it into a clean PDF. This file is a complete CV. All the
%  names, companies and numbers below are PLACEHOLDERS: replace them with yours.
%
%  HOW TO TURN THIS FILE INTO A PDF (pick one)
%    a) Overleaf (no install): overleaf.com -> New Project -> Upload Project ->
%       upload this .tex file -> click "Recompile". Compiler must be pdfLaTeX
%       (Menu -> Settings -> Compiler), which is Overleaf's default.
%    b) On your computer (needs TeX Live or MacTeX installed), in a terminal:
%           latexmk -pdf cv_template.tex
%       or, if you do not have latexmk:
%           pdflatex cv_template.tex
%
%  HOW THIS FILE IS ORGANISED
%    PART A - SETTINGS   (page size, fonts, spacing, layout helpers).
%             You can leave this part alone. Each number that is worth
%             tweaking has a comment saying what happens if you change it.
%    PART B - YOUR CONTENT (starts at \begin{document}). This is where you type.
%
%  THE 5 RULES THAT PREVENT 90% OF BEGINNER ERRORS
%    1. A % starts a COMMENT: everything after it on that line is ignored and
%       never appears in the PDF. (That is what all these lines are.)
%    2. Braces { } always come in pairs. Every { needs its matching }.
%    3. Every \begin{something} needs a matching \end{something}.
%    4. These characters are special. To print them, put a backslash first:
%           %  -> \%      &  -> \&      $  -> \$      _  -> \_      #  -> \#
%    5. A blank line starts a new paragraph. Single line breaks are ignored,
%       so you can split long sentences over several lines freely.
%
%  HANDY FORMATTING
%    \textbf{bold}   \textit{italic}   \textit{\textbf{bold italic}}
%    --  (two hyphens) gives a date-range dash: JAN 2024 -- MAR 2025
%    $\rightarrow$   gives an arrow:  30 $\rightarrow$ 360
%    \vspace{0.2 cm} adds vertical space; a negative value pulls things up.
%
%  WHEN SOMETHING BREAKS
%    The compiler reports the LINE NUMBER of the first error. Look there and
%    just above it for: a missing }, an unescaped % or &, or a \begin without
%    its \end. Fix the FIRST error first; later ones are often side effects.
% =============================================================================


% =============================================================================
%  PART A - SETTINGS  (safe to leave as is)
% =============================================================================

% -----------------------------------------------------------------------------
% A1. Document type, base font size and paper size
% -----------------------------------------------------------------------------
% 10pt  = body text size. Allowed: 10pt, 11pt, 12pt (bigger = longer CV).
% letterpaper = US Letter. Use a4paper for A4 (common outside North America).
\documentclass[10pt, letterpaper]{article}

% -----------------------------------------------------------------------------
% A2. Page margins
% -----------------------------------------------------------------------------
% Smaller margins fit more text on a page. 1 cm left/right is already tight;
% going lower may get cut off by some printers.
\usepackage[
    ignoreheadfoot,  % margins ignore header/footer space (this CV has neither)
    top=2 cm,        % space between the page's top edge and the text
    bottom=2 cm,     % space between the page's bottom edge and the text
    left=1 cm,       % space between the page's left edge and the text
    right=1 cm,      % space between the page's right edge and the text
    footskip=1.0 cm, % distance to the footer (unused, kept for completeness)
    % showframe,     % remove the % at the start of this line to draw the
                     % margin boxes - useful when adjusting layout
]{geometry}

% -----------------------------------------------------------------------------
% A3. Packages (add-ons). Each line loads one tool this CV uses.
% -----------------------------------------------------------------------------
\usepackage{titlesec}            % lets us restyle the \section headings
\usepackage[dvipsnames]{xcolor}  % colours
\usepackage{enumitem}            % lets us restyle bullet lists
\usepackage{fontawesome5}        % icons: \faEnvelope, \faPhone, \faLinkedin ...
\usepackage{changepage}          % 'adjustwidth', used by the entry blocks
\usepackage{paracol}             % side-by-side columns (text left, date right)
\usepackage{needspace}           % keeps a heading from being stranded at the
                                 % bottom of a page without its content

% Main colour used for links. Black means links look like normal text, which
% is what most recruiters expect. Try {RGB}{0, 79, 144} for a dark blue.
\definecolor{primaryColor}{RGB}{0, 0, 0}

% Clickable links + PDF properties (the "File > Properties" info of the PDF).
% Change pdftitle and pdfauthor to your own name.
\usepackage[
    pdftitle={Alex Morgan CV},
    pdfauthor={Alex Morgan},
    colorlinks=true,          % colour the link text instead of drawing boxes
    urlcolor=primaryColor,    % colour of web/email links
]{hyperref}

% -----------------------------------------------------------------------------
% A4. Fonts and machine-readable text (helps Applicant Tracking Systems)
% -----------------------------------------------------------------------------
% With pdfLaTeX, these lines make the PDF's text copy-pasteable and searchable,
% so automated CV scanners can read it. They do not change how it looks.
\usepackage{iftex}
\ifPDFTeX
    \input{glyphtounicode}
    \pdfgentounicode=1
    \usepackage[T1]{fontenc}
\fi

% Body font: Charter (a clean, readable serif). Delete this line to get the
% default LaTeX font, or try \usepackage{helvet} with
% \renewcommand{\familydefault}{\sfdefault} for a sans-serif look.
\usepackage{charter}

% -----------------------------------------------------------------------------
% A5. General page behaviour
% -----------------------------------------------------------------------------
\pagestyle{empty}                 % no header or footer
\pagenumbering{gobble}            % no page numbers
\setcounter{secnumdepth}{0}       % headings are not numbered (no "1 Education")
\setlength{\parindent}{0pt}       % paragraphs are not indented
\setlength{\topskip}{0pt}         % no extra space at the top of each page
\setlength{\columnsep}{0.15cm}    % gap between the text column and date column
\AtBeginEnvironment{adjustwidth}{\partopsep0pt} % no extra space before entries

% -----------------------------------------------------------------------------
% A6. Section heading style  (e.g. "EXPERIENCE" with a line under it)
% -----------------------------------------------------------------------------
% \needspace{4\baselineskip}: if fewer than ~4 lines are left on the page,
%     start the heading on the next page instead.
% \bfseries\large: bold and slightly bigger. \titlerule: the horizontal line.
\titleformat{\section}
    {\needspace{4\baselineskip}\bfseries\large}  % look of the heading
    {}{0pt}{}                                    % (no numbers)
    [\vspace{1pt}\titlerule]                     % line under the heading

% Space around every heading: {left}{above}{below}.
% Increase "above" to separate sections more; decrease to save space.
\titlespacing{\section}{-1pt}{0.3 cm}{0.2 cm}

% -----------------------------------------------------------------------------
% A7. Layout building blocks used in PART B
% -----------------------------------------------------------------------------
% You will use three blocks over and over:
%
%   onecolentry  - a full-width block of text.
%       \begin{onecolentry}
%           Any text here.
%       \end{onecolentry}
%
%   twocolentry  - text on the left, short text (usually dates) on the right.
%       \begin{twocolentry}{RIGHT SIDE, e.g. dates}
%           LEFT SIDE, e.g. company and job title
%       \end{twocolentry}
%       (Yes, the RIGHT side is written first, in the { } after \begin.)
%
%   highlights   - a bullet list. One \item per bullet.
%       \begin{highlights}
%           \item First bullet
%           \item Second bullet
%       \end{highlights}

% Bullet symbol: a small, vertically centred dot.
\renewcommand\labelitemi{$\vcenter{\hbox{\small$\bullet$}}$}

% highlights = a compact bullet list.
\newenvironment{highlights}{
    \begin{itemize}[
        topsep=0.10 cm,   % space above and below the whole list
        parsep=0.10 cm,   % space between bullets (bigger = airier list)
        partopsep=0pt,
        itemsep=0pt,
        leftmargin=10pt   % how far bullets are indented from the left
    ]
}{
    \end{itemize}
}

% onecolentry = a full-width block (practically no extra indent left/right).
% IMPORTANT: keep 0.00001 cm. It looks pointless, but with exactly 0 cm the
% date column of every twocolentry below silently disappears from the PDF.
\newenvironment{onecolentry}{
    \begin{adjustwidth}{0.00001 cm}{0.00001 cm}
}{
    \end{adjustwidth}
}

% twocolentry = left column takes the remaining width, right column is 4.5 cm.
% If your dates wrap onto two lines, increase 4.5 cm; if company names wrap
% too early, decrease it. This one number affects EVERY two-column entry.
% Tip: keep the left text to one line. If it wraps, an empty gap appears
% under the dates - shorten the text (e.g. "DEPT." instead of "DEPARTMENT").
\newenvironment{twocolentry}[1]{
    \onecolentry
    \def\secondColumn{#1}
    \setcolumnwidth{\fill, 4.5 cm}
    \begin{paracol}{2}
}{
    \switchcolumn \raggedleft \secondColumn
    \end{paracol}
    \endonecolentry
}

% header = centred block for your name and contact details.
% \linespread{1.5} = line spacing inside the header.
\newenvironment{header}{
    \setlength{\topsep}{0pt}\par\kern\topsep\centering\linespread{1.5}
}{
    \par\kern\topsep
}

% \contactitem{ICON}{LINK}{TEXT SHOWN}  - one clickable contact detail.
%   ICON: an icon command such as \faEnvelope (find more icons by searching
%         "fontawesome5 package documentation").
%   LINK: where a click goes: mailto:you@x.com, tel:+15550100, https://...
%   TEXT: what is printed. Keep LINK and TEXT in sync when you edit them!
\newcommand{\contactitem}[3]{\mbox{\href{#2}{#1\hspace{1pt} #3}}}

% \contactsep  - the " | " separator between contact items.
% (\AND draws the bar; \unskip and \ignorespaces remove stray spaces.)
\newsavebox\ANDbox
\sbox\ANDbox{$|$}
\newcommand{\AND}{\unskip\cleaders\copy\ANDbox\hskip\wd\ANDbox\ignorespaces}
\newcommand{\contactsep}{\kern 5pt\AND\kern 5pt}


% =============================================================================
%  PART B - YOUR CONTENT  (edit everything below)
% =============================================================================
% Sections appear in the PDF in the same order as they appear here.
% To remove a section: delete it, or put % in front of each of its lines.
% To reorder: cut and paste whole sections (from one \section to the next).
% To rename a heading: change the text inside \section{...}.

\begin{document}

    % -------------------------------------------------------------------------
    % NAME AND CONTACT DETAILS
    % -------------------------------------------------------------------------
    % \fontsize{25 pt}{25 pt}: first number = name size, second = line height.
    % Delete a \contactitem line (and the \contactsep before it) to drop it.
    \begin{header}
        \fontsize{25 pt}{25 pt}\selectfont Alex Morgan

        \vspace{5 pt}

        \normalsize
        \contactitem{\faEnvelope}{mailto:alex.morgan@example.com}{alex.morgan@example.com}\contactsep
        \contactitem{\faPhoneSquare}{tel:+15550100}{+1 555 0100}\contactsep
        \contactitem{\faLinkedin}{https://www.linkedin.com/in/your-profile}{linkedin.com/in/your-profile}\contactsep
        \contactitem{\faGlobe}{https://example.com}{example.com}
    \end{header}

    % Pulls the first section slightly closer to the contact line.
    \vspace{-0.12 cm}

    % -------------------------------------------------------------------------
    % PROFESSIONAL SUMMARY
    % -------------------------------------------------------------------------
    % 2-4 sentences: who you are, years of experience, your biggest results,
    % and what you bring. Write numbers as \$300M (the $ must be escaped).
    \section{PROFESSIONAL SUMMARY}
    \begin{onecolentry}
        Senior data scientist and applied AI engineer with 3+ years of industry experience across cloud software and consumer technology at Google and Microsoft. Built interpretable forecasting models and scaled personalization systems that contributed over \$300M in incremental revenue, and shipped LLM-powered agent workflows that automate internal operations, backed by a combined bachelor's and master's degree in Electrical Engineering.
    \end{onecolentry}

    % -------------------------------------------------------------------------
    % EDUCATION
    % -------------------------------------------------------------------------
    % Two rows: (1) institution + dates, (2) degree + GPA.
    % Copy both twocolentry blocks to add another degree.
    \section{EDUCATION}

    \begin{twocolentry}{
        \textbf{AUG 2017 -- MAY 2022}
        }
        \textbf{THE UNIVERSITY OF TEXAS AT AUSTIN}
    \end{twocolentry}
    \vspace{0.05 cm}

    \begin{twocolentry}{
        \textit{GPA: 3.7/4.0}
        }
        \textit{B.S. + M.S. in Electrical Engineering}
    \end{twocolentry}

    % -------------------------------------------------------------------------
    % EXPERIENCE
    % -------------------------------------------------------------------------
    % RECIPE - one job = a twocolentry (title row) + a onecolentry (bullets):
    %
    %   \begin{twocolentry}{
    %       \textbf{START -- END}
    %       }
    %       \textbf{COMPANY} \textit{(CITY, STATE)}, JOB TITLE
    %   \end{twocolentry}
    %   \begin{onecolentry}
    %       \vspace{0.05 cm}
    %       \begin{highlights}
    %           \item What you did, how, and the measurable result
    %       \end{highlights}
    %   \end{onecolentry}
    %   \vspace{0.2 cm}      <- gap before the next job
    %
    % Bullet tip: start with a strong verb (Built, Led, Reduced...), and end with
    % a number (\%, \$, time saved). Bold only the 2-3 words you want skimmed.
    \section{EXPERIENCE}

    % Job 1: a simple job with one list of bullets.
    \begin{twocolentry}{
        \textbf{JAN 2025 -- PRESENT}
        }
        \textbf{GOOGLE} \textit{(MOUNTAIN VIEW, CA)}, SENIOR AI ENGINEER
    \end{twocolentry}
    \begin{onecolentry}
        \vspace{0.05 cm}
        \begin{highlights}
            \item Redesigned a manual, multi-team device-return process into a single \textbf{auditable pipeline}, consolidating ticket validation, inventory updates and refund approvals behind one \textbf{governed execution path}
            \item Built a \textbf{Kotlin/Spring Boot service} exposing REST APIs over parameterized PostgreSQL queries, using OAuth 2.0 and \textbf{Secret Manager-backed configuration} to replace brittle legacy scripts without changing business rules
            \item Stabilized \textbf{LLM outputs} across multi-step tasks with \textbf{JSON schemas}, \textbf{input validation} and curated \textbf{in-context examples}, benchmarking sampling settings across several \textbf{commercial and open-weight models}
            \item Launched a \textbf{human-in-the-loop agent workflow} that handles requests in parallel and produces one summary report for operators, cutting end-to-end \textbf{handling time by 85\%}
        \end{highlights}
    \end{onecolentry}
    \vspace{0.2 cm}

    % Job 2: a job split into several project groups, each with its own title
    % (\textit{\textbf{...}}) and bullet list. 0.15 cm separates the groups.
    \begin{twocolentry}{
        \textbf{JUL 2022 -- DEC 2024}
        }
        \textbf{MICROSOFT} \textit{(REDMOND, WA)}, SENIOR DATA SCIENTIST
    \end{twocolentry}
    \begin{onecolentry}
        \vspace{0.1 cm}
        \textit{\textbf{Customer Analytics}}
    \end{onecolentry}
    \begin{onecolentry}
        \begin{highlights}
            \item Built a \textbf{LightGBM}-based \textbf{subscription usage forecasting} model, applying feature selection to cut the input feature set by 70\% while speeding up training and improving \textbf{interpretability}
            \item Strengthened model reliability with \textbf{Bayesian hyperparameter search} and \textbf{SHAP-based explanations}, turning an opaque model into a reviewable system that met internal \textbf{responsible-AI \& compliance} standards
            \item Lowered forecast error by \textbf{18\%} (MAPE), adding an estimated \textbf{USD 15M} in \textbf{annual incremental revenue} and extending the \textbf{retraining interval 6x} (monthly $\rightarrow$ semi-annual), which cut maintenance costs
        \end{highlights}
    \end{onecolentry}
    \vspace{0.15 cm}
    \begin{onecolentry}
        \textit{\textbf{Marketing Data Platform}}
    \end{onecolentry}
    \begin{onecolentry}
        \begin{highlights}
            \item Led the build-out of a shared \textbf{feature store} for marketing analytics, letting 25+ downstream models \textbf{reuse common features} and providing a \textbf{scalable data foundation} for new use cases
            \item Applied time-window and rule-based de-duplication to \textbf{shrink the feature catalog} by 35\%, and defined a \textbf{single source of truth} for every feature, shortening model development cycles
            \item Drove the \textbf{migration of training pipelines} from on-premise clusters to \textbf{Microsoft Azure}, with data reconciliation, logic standardization and automated regression tests to guarantee \textbf{no data loss}
        \end{highlights}
    \end{onecolentry}
    \vspace{0.15 cm}
    \begin{onecolentry}
        \textit{\textbf{Pricing and Promotions}}
    \end{onecolentry}
    \begin{onecolentry}
        \begin{highlights}
            \item Projected 12-month \textbf{revenue} and \textbf{cost} for a customer rewards promotion by comparing forecast and actual redemption history, informing \textbf{budget planning} and quarterly \textbf{financial targets}
            \item Analyzed past \textbf{promotion redeemers} to derive \textbf{calibration factors}, refine the incentive framework and surface \textbf{high-ROI} customer cohorts with \textbf{shifting purchase behavior}
            \item Fed calibrated forecasts into the \textbf{personalized offers} engine, improving funding decisions and raising \textbf{incremental revenue} from USD 200M to \textbf{USD 310M} while protecting margins
        \end{highlights}
    \end{onecolentry}

    % -------------------------------------------------------------------------
    % AWARDS AND ACHIEVEMENTS
    % -------------------------------------------------------------------------
    % A plain bullet list (no onecolentry needed). One \item per award.
    \section{AWARDS AND ACHIEVEMENTS}
    \begin{highlights}
        \item \textbf{Microsoft Team Excellence Award -- 2024 (Impact)} for leading the marketing data platform initiative
        \item \textbf{First place at a national collegiate robotics challenge} as team lead for an autonomous rover design event
    \end{highlights}

    % -------------------------------------------------------------------------
    % INTERNSHIPS
    % -------------------------------------------------------------------------
    % Same recipe as EXPERIENCE. The location is optional - here it is omitted.
    \section{INTERNSHIPS}
    \begin{twocolentry}{
        \textbf{MAY 2021 -- AUG 2021}
        }
        \textbf{AMAZON}, SOFTWARE DEVELOPMENT ENGINEER INTERN
    \end{twocolentry}
    \begin{onecolentry}
        \vspace{0.05 cm}
        \begin{highlights}
            \item Built an internal web portal for \textbf{equipment booking and tracking} used by more than 30K employees
            \item Designed reusable \textbf{back-end \& front-end modules} for reading and writing a \textbf{relational database} through web forms
            \item Tested 10+ modules across 60+ \textbf{scenarios}, adding \textbf{clear validation messages} that guided users to fix input errors
        \end{highlights}
    \end{onecolentry}
    \vspace{0.2 cm}
    \begin{twocolentry}{
        \textbf{JUN 2019 -- AUG 2019}
        }
        \textbf{SPOTIFY}, DATA SCIENCE INTERN
    \end{twocolentry}
    \begin{onecolentry}
        \vspace{0.05 cm}
        \begin{highlights}
            \item Built a personalized \textbf{content recommendation engine} for a music discovery feature to increase listener engagement
            \item Trained \textbf{item embeddings} from historical listening sessions with a \textbf{skip-gram model} implemented in TensorFlow
            \item Surfaced recommendations in the app, lifting average \textbf{listening time} by 30\% and \textbf{track saves} by 45\%
        \end{highlights}
    \end{onecolentry}

    % -------------------------------------------------------------------------
    % PROJECTS
    % -------------------------------------------------------------------------
    % Same recipe as EXPERIENCE: PROJECT NAME, lab/department on the left,
    % dates on the right, 2-3 bullets underneath.
    \section{PROJECTS}
    \begin{twocolentry}{
        \textbf{JAN 2022 -- MAY 2022}
        }
        \textbf{PREDICTIVE MAINTENANCE PLATFORM}, SMART MANUFACTURING LAB, UT AUSTIN
    \end{twocolentry}
    \begin{onecolentry}
        \vspace{0.05 cm}
        \begin{highlights}
            \item Created a \textbf{real-time simulation} of an industrial motor rig to support \textbf{anomaly detection} and \textbf{remaining-life estimation}
            \item Built a \textbf{sensor data pipeline} covering acquisition, \textbf{noise filtering} and storage, with automated \textbf{feature extraction}
            \item Trained a \textbf{classification model} that flags abnormal vibration patterns early to enable \textbf{preventive maintenance}
        \end{highlights}
    \end{onecolentry}
    \vspace{0.2 cm}
    \begin{twocolentry}{
        \textbf{AUG 2021 -- DEC 2021}
        }
        \textbf{FAIRNESS IN SOCIAL RECOMMENDATIONS}, DEPT. OF COMPUTER SCIENCE, UT AUSTIN
    \end{twocolentry}
    \begin{onecolentry}
        \vspace{0.05 cm}
        \begin{highlights}
            \item Studied how \textbf{algorithmic bias in friend suggestions} can create echo chambers on social platforms, measuring group separation with \textbf{graph centrality metrics} and within- versus across-group distances
            \item Designed a simulation of \textbf{user activity} and \textbf{network growth} based on a neutral probabilistic model
            \item Built a \textbf{recommendation model} using \textbf{graph neural networks} trained on \textbf{link prediction} tasks
        \end{highlights}
    \end{onecolentry}
    \vspace{0.2 cm}
    \begin{twocolentry}{
        \textbf{JAN 2021 -- MAY 2021}
        }
        \textbf{INTERPRETABLE TEXT EMBEDDINGS}, NATURAL LANGUAGE PROCESSING LAB, UT AUSTIN
    \end{twocolentry}
    \begin{onecolentry}
        \vspace{0.05 cm}
        \begin{highlights}
            \item Implemented a \textbf{sense-aware word embedding} model that splits ambiguous words into their distinct meanings
            \item Benchmarked it against \textbf{GloVe} embeddings on similarity and inference tasks, improving \textbf{F1 score} by up to 12\%
            \item Combined \textbf{embeddings} with \textbf{NMF topic modeling} to produce sparse, readable topic vectors with strong \textbf{coherence scores}
        \end{highlights}
    \end{onecolentry}

    % -------------------------------------------------------------------------
    % SKILLS AND EXPERTISE
    % -------------------------------------------------------------------------
    % One onecolentry per skill group: \textbf{Group name:} then a comma list.
    % Keep each group to about one line. 0.05 cm = small gap between groups.
    \section{SKILLS AND EXPERTISE}
    \begin{onecolentry}
        \textbf{Programming and Data:} Python (Pandas, Polars, scikit-learn, spaCy), Kotlin, SQL, Snowflake, Apache Spark, Airflow, dbt
    \end{onecolentry}
    \vspace{0.05 cm}
    \begin{onecolentry}
        \textbf{Machine Learning:} LightGBM, CatBoost, SHAP, Bayesian optimization, time-series forecasting, TensorFlow, causal inference
    \end{onecolentry}
    \vspace{0.05 cm}
    \begin{onecolentry}
        \textbf{NLP:} text embeddings, NMF topic modeling, named-entity recognition, semantic similarity evaluation
    \end{onecolentry}
    \vspace{0.05 cm}
    \begin{onecolentry}
        \textbf{Generative AI and Agents:} LLM application development, prompt design, structured outputs, RAG, tool calling
    \end{onecolentry}
    \vspace{0.05 cm}
    \begin{onecolentry}
        \textbf{Cloud and Tools:} Microsoft Azure, Google Cloud, Docker, REST APIs, Git, Jupyter, VS Code
    \end{onecolentry}
    \vspace{0.05 cm}
    \begin{onecolentry}
        \textbf{Business and Delivery:} A/B testing, ROI analysis, revenue forecasting, Scrum, Jira, stakeholder management
    \end{onecolentry}

    % -------------------------------------------------------------------------
    % CERTIFICATIONS
    % -------------------------------------------------------------------------
    % \item \textbf{Certificate name}, Issuer
    \section{CERTIFICATIONS}
    \begin{highlights}
        \item \textbf{Azure AI Engineer Associate and Azure Data Scientist Associate}, Microsoft
        \item \textbf{Professional Machine Learning Engineer}, Google Cloud
        \item \textbf{Machine Learning Engineer -- Associate}, Amazon Web Services
        \item \textbf{Deep Learning Specialization}, DeepLearning.AI
    \end{highlights}

    % -------------------------------------------------------------------------
    % EXTRACURRICULAR ACTIVITIES
    % -------------------------------------------------------------------------
    % \textbf{Role or activity}: one line on what you do and why it matters.
    \section{EXTRACURRICULAR ACTIVITIES}
    \begin{onecolentry}
        \textbf{Volunteer}: Organizing weekend food drives with a local food bank to support families facing food insecurity
    \end{onecolentry}
    \vspace{0.05 cm}
    \begin{onecolentry}
        \textbf{Mentor, University Alumni Network}: Advising students on career paths, interview preparation and learning plans
    \end{onecolentry}
    \vspace{0.05 cm}
    \begin{onecolentry}
        \textbf{Organizer, Internal Tech Talks}: Hosting monthly sessions where teams share their work and explore cross-team roles
    \end{onecolentry}
    \vspace{0.05 cm}
    \begin{onecolentry}
        \textbf{Events Lead}: Coordinating team offsites and celebrations that recognize milestones and build team spirit
    \end{onecolentry}
    \vspace{0.05 cm}
    \begin{onecolentry}
        \textbf{Hiking and Writing}: Exploring national parks and writing short travel essays about the people and places met along the way
    \end{onecolentry}

\end{document}
